\chapter{Aire minimale de la coupe triadique, capacité \texorpdfstring{\(81\log 3\)}{81 log 3} et transport collectif d’incidence}
\label{chap:parte-iv-58}

La quantification aréale part de l’intérieur nonadique, de son bord et du transport collectif d’incidence. La coupe minimale contient \(81\) unités et sa capacité qutrit est \(81\log 3\) ; le résultat précède sa réalisation gravitationnelle.

\section{Quantification aréale, dynamique modulaire et transformation rationnelle CKM}
\label{rm:ch:modular}

\subsection{Inclusion hexade--octade et réalisations fonctionnelles associées}

Fixons un ensemble \(H\) de six points et un ensemble \(O\) de huit
points avec une inclusion distinguée \(j:H\hookrightarrow O\). Dans la
résidence de Witt, \(H\) est une hexade d'une dodécade et \(O\) son octade
relevée ; dans ce chapitre, on utilise seulement l'inclusion déjà fixée, et non la
valeur terminale de Barbero--Immirzi. Nous écrivons
\begin{equation}
\label{rm:eq:modular-pairs-H}
\binom H2:=\{P\subset H:|P|=2\},
\qquad \left|\binom H2\right|=\binom62=15,
\end{equation}
et définissons deux espaces d'incidences de types distincts :
\begin{equation}
\label{rm:eq:modular-F6-F8}
\mathcal F_6:=H\times\binom H2,
\qquad
\mathcal F_8:=O\times\binom H2.
\end{equation}
L'inclusion d'incidences est
\begin{equation}
\label{rm:eq:modular-iota68}
\iota_{6,8}:\mathcal F_6\hookrightarrow\mathcal F_8,
\qquad
(h,P)\longmapsto(j(h),P).
\end{equation}
Elle est injective parce que \(j\) l'est et conserve la paire marquée \(P\).

\begin{theorem}[Décomptes d'incidence et défaut orienté]
\label{rm:thm:modular-incidence-90-120}
Les deux espaces de \cref{rm:eq:modular-F6-F8} satisfont
\begin{equation}
\label{rm:eq:modular-counts-90-120}
|\mathcal F_6|=6\binom62=90,
\qquad
|\mathcal F_8|=8\binom62=120,
\qquad
|\mathcal F_8\setminus\iota_{6,8}(\mathcal F_6)|=30.
\end{equation}
La normalisation par les vingt-quatre coordonnées du couple
dodécade--vingt-quatre et par les quinze paires de \(H\) produit
\begin{equation}
\label{rm:eq:modular-incidence-defect}
\boxed{
\delta_{6,8}^{\rm inc}
=\frac{|\mathcal F_6|-|\mathcal F_8|}
{24\binom62}=-\frac1{12}.}
\end{equation}
La même fraction s'obtient par la fonctionnelle réciproque associée
à la mémoire incorporée,
\begin{equation}
\label{rm:eq:modular-reciprocal-defect}
\boxed{
\delta_{6,8}^{\rm rec}
=\frac{30}{120}-\frac{30}{90}=-\frac1{12}.}
\end{equation}
Ce sont deux réalisations exactes du même défaut orienté, mais elles ont
des domaines distincts et ne sont pas identifiées à \(\zeta(-1)\) sans un
entrelaceur supplémentaire.
\end{theorem}

\begin{proof}
Le principe multiplicatif appliqué à
\cref{rm:eq:modular-F6-F8} donne \(6\cdot15=90\) et \(8\cdot15=120\).
L'image de \(\iota_{6,8}\) a quatre-vingt-dix éléments, de sorte que le
complémentaire en a trente. Substituer dans
\cref{rm:eq:modular-incidence-defect} donne
\(-30/(24\cdot15)=-1/12\), tandis que
\cref{rm:eq:modular-reciprocal-defect} est
\(1/4-1/3=-1/12\).
\end{proof}

La transition hexade--octade admet deux réalisations fonctionnelles qui ne
se déduisent pas l'une de l'autre. Le diagramme algébrique, formulé sur les espaces
qui justifient les décomptes, est
\begin{equation}
\label{rm:eq:modular-68-two-readers}
\begin{aligned}
(H\xhookrightarrow{\,j\,}O)
&\longmapsto
\left(|\mathcal F_6|,|\mathcal F_8|\right)=(90,120)
\longmapsto\Gamma_{6\to8}=\gammaH,\\
(h_6,o_8,t)=(6,8,3)
&\longmapsto-\frac{o_8-h_6}{t}=-\frac23,
\end{aligned}
\end{equation}
où \(-2/3\) est le coefficient de la troisième coordonnée de la ligne
\(13\) du constructeur CKM. La première réalisation transforme
l'incidence en aire et en mémoire modulaire ; la seconde transforme la même
transition en paramètres de mélange angulaire.
On ne postule pas \(\boldsymbol\theta_{\rm CKM}=f(\gammaH)\) : les deux sorties
conservent un antécédent commun.

Le morphisme qui détermine les secteurs \(90/120\) est \(\iota_{6,8}\). La grandeur \(\gammaH\) est adoptée comme théorème déjà démontré dans la construction hexade--octade et n’est pas reconstruite à partir de sa représentation décimale. Les conséquences opératorielles ne commencent qu’après fixation de l’échelle élémentaire \(a_0\), définie ci-dessous sans utiliser l’entropie ni le produit terminal \(\gammaH a_0\).

\subsection{Correspondance entre l'intérieur nonadique et les degrés de liberté de bord}

Soit \(\mathbb T_3=\{0,1,2\}\). Tout \(x\in\mathbb Z/9\mathbb Z\) possède
une écriture en chiffres unique
\(x=x_0+3x_1\), avec \(x_i\in\mathbb T_3\), et tout
\(y\in\mathbb Z/27\mathbb Z\) une écriture unique
\(y=y_0+3y_1+9y_2\). C'est pourquoi
\begin{equation}
\label{rm:eq:modular-bulk-boundary-729}
\mathcal B=(\mathbb Z/9\mathbb Z)^3,
\qquad
\mathcal B_\partial=(\mathbb Z/27\mathbb Z)^2,
\qquad
|\mathcal B|=|\mathcal B_\partial|=3^6=729.
\end{equation}
L'égalité des cardinaux s'accompagne de l'application explicite. Si
\(x_j=r_{2j-1}+3r_{2j}\) pour \(j=1,2,3\), nous définissons
\begin{equation}
\label{rm:eq:modular-beta729}
\beta_{729}(x_1,x_2,x_3)
=\bigl(r_1+3r_2+9r_3,\;r_4+3r_5+9r_6\bigr).
\end{equation}

\begin{proposition}[Bijection généalogique \(729\leftrightarrow729\)]
\label{rm:prop:modular-beta729}
L'application \(\beta_{729}:\mathcal B\to\mathcal B_\partial\) est bijective et
préserve le mot ordonné de six trits. Ce n'est pas, en général, un
isomorphisme de groupes additifs : regrouper les chiffres change l'endroit où
la retenue est enregistrée.
\end{proposition}

\begin{proof}
Extraire les six chiffres de trois coordonnées nonadiques et les regrouper en
deux triplets est une permutation des coordonnées de \(\mathbb T_3^6\). La
décomposition positionnelle est unique des deux côtés, de sorte que
l'opération possède un inverse. La mise en garde additive se vérifie, par exemple,
en additionnant des mots qui engendrent une retenue entre \(r_3\) et \(r_4\) : cette
frontière sépare les deux coordonnées d'ordre vingt-sept, mais traverse
deux coordonnées distinctes dans le regroupement nonadique.
\end{proof}

Sur le tore de bord \(\mathcal B_\partial\), nous prenons le bloc de sommets
\begin{equation}
\label{rm:eq:modular-block-A}
A=\{0,1,\ldots,8\}\times\{0,1,\ldots,8\}.
\end{equation}
Ses liens orientés qui traversent le bord se séparent par direction :
\begin{align}
\partial A_{90}
&=\{((-1,j),(0,j)),((8,j),(9,j)):0\leq j\leq8\},
\label{rm:eq:modular-boundary90}\\
\partial A_{120}
&=\{((i,-1),(i,0)),((i,8),(i,9)):0\leq i\leq8\},
\label{rm:eq:modular-boundary120}
\end{align}
où toutes les coordonnées se lisent modulo vingt-sept. Ainsi
\begin{equation}
\label{rm:eq:modular-boundary-counts}
\partial A=\partial A_{90}\sqcup\partial A_{120},
\qquad
|\partial A_{90}|=|\partial A_{120}|=18.
\end{equation}
Les indices \(90,120\) étiquettent la provenance du canal ; ils ne disent pas que
le bord possède quatre-vingt-dix ou cent vingt liens.

\begin{theorem}[Loi aréale de la coupe triadique]
\label{rm:thm:modular-triadic-cut}
Dans le graphe cubique \(N\times N\times N\) à capacité unitaire, toute
coupe entre deux faces opposées possède une capacité d'au moins \(N^2\), et une
couche transversale atteint cette borne. Pour \(N=9\),
\begin{equation}
\label{rm:eq:modular-mincut-entropy}
\mathcal A_{\min}=81,
\qquad
S_{\rm tri}=81\log3.
\end{equation}
\end{theorem}

\begin{proof}
Il existe \(N^2\) droites d'arêtes parallèles à la direction normale, deux à deux
disjointes. Toute coupe doit intercepter chacune d'elles, de sorte que sa capacité
est d'au moins \(N^2\). Une seule couche transversale contient exactement
\(N^2\) arêtes et réalise le minimum. Si chaque lien coupé porte un trit
maximalement mélangé, son entropie est \(\log3\) ; l'additivité sur les
quatre-vingt-un liens donne \cref{rm:eq:modular-mincut-entropy}.
\end{proof}

L'entropie \(S_{\rm tri}\) mesure la capacité de la coupe cubique.
L'entropie modulaire des trente-six liens de
\cref{rm:eq:modular-boundary-counts} mesure un autre état et ne doit pas lui être
égalisée.

\subsection{Entrelaceur collectif d'incidence entre intérieur et bord}

Dans la somme
\(\ell^2(\mathcal F_6)\oplus\ell^2(\mathcal F_8)\), nous définissons
\begin{equation}
\label{rm:eq:modular-u90-u120}
u_{90}=\frac1{\sqrt{90}}\sum_{f\in\mathcal F_6}(\delta_f,0),
\qquad
u_{120}=\frac1{\sqrt{120}}\sum_{f\in\mathcal F_8}(0,\delta_f).
\end{equation}
Dans
\(\ell^2(\partial A_{90})\oplus\ell^2(\partial A_{120})\), soient
\begin{equation}
\label{rm:eq:modular-v90-v120}
v_{90}=\frac1{\sqrt{18}}\sum_{e\in\partial A_{90}}(\delta_e,0),
\qquad
v_{120}=\frac1{\sqrt{18}}\sum_{e\in\partial A_{120}}(0,\delta_e).
\end{equation}
Les couples \((u_{90},u_{120})\) et \((v_{90},v_{120})\) sont orthonormaux.
Il existe donc une unique isométrie entre les sous-espaces collectifs
qui satisfait
\begin{equation}
\label{rm:eq:modular-J68}
\mathcal J_{6\to8}u_{90}=v_{90},
\qquad
\mathcal J_{6\to8}u_{120}=v_{120}.
\end{equation}
Définissons
\begin{align}
D_{\rm inc}
&=90|u_{90}\rangle\langle u_{90}|
+120|u_{120}\rangle\langle u_{120}|,
\label{rm:eq:modular-Dinc}\\
D_{\rm area}
&=90|v_{90}\rangle\langle v_{90}|
+120|v_{120}\rangle\langle v_{120}|.
\label{rm:eq:modular-Darea}
\end{align}

\begin{theorem}[Transport exact des deux étiquettes]
\label{rm:thm:modular-J68}
L'application \(\mathcal J_{6\to8}\) est unitaire entre les deux plans
collectifs et entrelace les opérateurs d'incidence et d'aire :
\begin{equation}
\label{rm:eq:modular-J-intertwining}
\boxed{
\mathcal J_{6\to8}D_{\rm inc}
=D_{\rm area}\mathcal J_{6\to8}.}
\end{equation}
Par conséquent, toute combinaison polynomiale des étiquettes \(90,120\), en
particulier la combinaison primitive \(12:-1\), est transportée avec les
mêmes coefficients.
\end{theorem}

\begin{proof}
L'orthonormalité montre que \(\mathcal J_{6\to8}\) envoie une base
orthonormale sur une autre. En \(u_{90}\), les deux membres de
\cref{rm:eq:modular-J-intertwining} valent \(90v_{90}\) ; en \(u_{120}\),
ils valent tous deux \(120v_{120}\). L'égalité sur une base prouve l'identité.
Le calcul fonctionnel polynomial conserve l'entrelacement.
\end{proof}

Cette application ne prétend pas injecter quatre-vingt-dix incidences individuelles dans
dix-huit liens : elle transporte exactement les deux modes uniformes et leurs
étiquettes spectrales. L'étendre à des fluctuations non uniformes exigerait
une nouvelle application d'incidence ; l'opérateur modulaire défini ici n'utilise que
le sous-espace collectif qui vient d'être construit.

\subsection{Quantum élémentaire et échelle minimale de l’opérateur d’aire}

Soit \(C^*\) la coordonnée angulaire conjuguée, exprimée en degrés, et fixons
\begin{equation}
\label{rm:eq:modular-theta-clock-a0}
\theta_{\rm clk}=\frac{C^*}{6}\frac{\pi}{180},
\qquad
\boxed{
a_0=\frac{1+2\cos(10^\circ)}{3}\,\theta_{\rm clk}.}
\end{equation}
Le facteur qui précède \(\theta_{\rm clk}\) n'est pas une correction
décimale. C'est le caractère réel normalisé du triplet de phases
\begin{equation}
\label{rm:eq:modular-c10-character}
\frac13\Tr\operatorname{diag}
\left(1,\ee^{\ii\pi/18},\ee^{-\ii\pi/18}\right)
=\frac{1+2\cos(10^\circ)}3.
\end{equation}
Ainsi, \(a_0\) est la moyenne spectrale de la structure angulaire sur le trit
orienté. Elle est positive, est fixée avant l'opérateur d'aire et n'est pas
sélectionnée à partir de l'entropie. L'identification ultérieure
\(\gammaH=\Gamma_{6\to8}\) réalise la sortie HMT dans la connexion
d'Ashtekar--Barbero \cite{Barbero1995,Immirzi1997} ; elle ne rouvre pas sa dérivation.

\begin{remark}[Portée du symbole \(a_0\)]
Dans ce chapitre, \(a_0\equiv a_0^{\rm area}\) désigne exclusivement la
normalisation de \cref{rm:eq:modular-theta-clock-a0}. Ce n'est ni le rayon de Bohr
ni le facteur d'échelle cosmologique initial, bien que ces objets s'écrivent
traditionnellement avec le même symbole. Aucune identité de ce chapitre
ne permet de substituer l'un à l'autre.
\end{remark}


\paragraph{Construction de l’opérateur d’aire à partir de la coupe triadique et séparation des canaux.} La coupe et sa capacité déterminent un opérateur d’aire ; son spectre doit conserver séparément l’occupation totale et la provenance du canal.
